\chapter{Experiments and Results}
\label{ch:results}

The three stages answer the research question at increasing realism. Stage~1
shows the optimum is task-dependent; Stage~2 shows a learnable memory is
competitive and characterizes what each parameter does; Stage~3 tests the idea
on real LLMs over six long-context corpora.

\section{Stage 1 --- task-dependent \texorpdfstring{$\thetavec$}{theta} in grid worlds}
\label{sec:stage1}
An Evolution Strategy optimizes the three-dimensional $\thetavec$ separately on
each grid environment. Table~\ref{tab:stage1} reports the recovered vectors and
the success rate against the fixed-memory baseline.

\begin{table}[t]
\centering
\caption{Stage~1: learned $\thetavec=(\theta_{\mathrm{store}},
\theta_{\mathrm{entity}}, \theta_{\mathrm{temporal}})$ per environment, and
baseline vs.\ learned success. The vectors are visibly different per task.}
\label{tab:stage1}
\begin{tabular}{lccc}
\toprule
Environment & Learned $\thetavec$ & Baseline & Learned \\
\midrule
Key-Door          & $(0.116,\,0.000,\,0.819)$ & 17.5\% & 30.0\% \\
Goal-Room         & $(1.000,\,0.220,\,1.000)$ & 70.0\% & 80.0\% \\
MultiHopKeyDoor   & $(1.000,\,0.487,\,0.843)$ & \textbf{2.5\%} & \textbf{27.5\%} \\
\bottomrule
\end{tabular}
\end{table}

The headline is the \emph{difference} between the recovered vectors:
Key-Door discards most events ($\theta_{\mathrm{store}}{=}0.12$) but preserves
temporal order; Goal-Room stores everything. This is the first, cleanest
evidence that the optimal memory structure is task-dependent
(Figure~\ref{fig:stage1theta}). On the hardest task the learned rule lifts
success from $2.5\%$ to $27.5\%$; we report this as an absolute lift of
$\sim$25 points (the baseline is near-zero, so a percentage-ratio framing would
overstate it), on a single-seed proof of concept.

\paragraph{Which parameters each task actually exercises.} Reading
Table~\ref{tab:stage1} column-wise rather than row-wise is more informative than
the raw success gap, because it shows \emph{which} of the three coordinates each
task puts under load. Key-Door is the diagnostic case: its recovered
$\theta_{\mathrm{store}}{=}0.116$ is the lowest of the three environments while
its $\theta_{\mathrm{temporal}}{=}0.819$ is near-maximal, and its
$\theta_{\mathrm{entity}}$ collapses to exactly $0.000$. That combination is not
an arbitrary point in the cube --- it is a precise behavioural description of the
task. A task whose only relevant fact is ``a key was seen, then a door reached''
needs to retain \emph{order} but almost nothing else, so the optimizer drives
storage down (discard most events) and entity weighting to the floor (the events
are not distinguished by \emph{what} object they involve) while holding temporal
structure high. Goal-Room exercises the opposite corner:
$\theta_{\mathrm{store}}{=}1.000$ and $\theta_{\mathrm{temporal}}{=}1.000$ say
that when the environment is small and information-dense the cheapest correct
policy is simply to store everything and never forget, and indeed its baseline is
already high ($70\%$) because a fixed default is accidentally close to that
store-all optimum. MultiHopKeyDoor is the task that most exercises the
\emph{entity} coordinate ($\theta_{\mathrm{entity}}{=}0.487$, by far the largest
of the three) precisely because multi-hop credit assignment requires the memory
to tell one keyed sub-goal from another; it is also the task where the fixed
baseline is most catastrophically wrong ($2.5\%$), and it couples a high store
weight with a high temporal weight, i.e.\ it loads all three axes at once.

\paragraph{Why no fixed default can serve the three tasks.} The three recovered
vectors do not merely differ in magnitude; they \emph{disagree in direction} on
every axis that matters. Storage is near-zero for Key-Door but saturated for the
other two; entity weight is exactly zero for Key-Door, moderate for
MultiHopKeyDoor, and low-but-nonzero for Goal-Room; only the temporal coordinate
is consistently high. A single fixed $\thetavec$ is a single point, so it can sit
at most near one of these corners --- which is exactly what the baseline column
shows: the default happens to land close to Goal-Room's store-all optimum
($70\%\rightarrow80\%$, a small lift because little was left to gain) and far from
Key-Door's discard-most optimum and from MultiHopKeyDoor's
selective-plus-ordered optimum ($2.5\%\rightarrow27.5\%$, the largest lift
precisely because the default was pointing the wrong way). The lift is therefore
not uniform ``more tuning helps everywhere''; it is inversely proportional to how
close the fixed default already sat to each task's corner. That is the sharpest
possible form of the task-dependence claim: the tasks do not want a
\emph{stronger} setting of the same knobs, they want \emph{opposing} settings, so
no grid-tuned constant can be simultaneously correct for all three, and
per-task adaptation is a requirement of the problem rather than an optimization
luxury.

\begin{figure}[t]
  \centering
  \includegraphics[width=0.9\linewidth]{fig3_theta_trajectory.png}
  \caption{Stage~1: $\thetavec$ trajectories over Evolution-Strategy generations
  for three environments. The optimizer converges to visibly different vectors,
  which is the central Stage-1 claim.}
  \label{fig:stage1theta}
\end{figure}

\section{Stage 2 --- benchmark, ablation, transfer, sensitivity}
\label{sec:stage2}

\paragraph{Top-cluster performance.} On the hardest grid task
(MultiHopKeyDoor), the CMA-ES-tuned ten-dimensional \Vfour{} reaches the top
performance cluster (Figure~\ref{fig:stage2bench}): mean reward $0.178$ versus
the prior best EpisodicSemantic at $0.173$. We state this as a
\emph{statistical tie}, not a clean ``\#1'': it is a single-seed point estimate
and EpisodicSemantic's $95\%$ bootstrap confidence interval $[0.120, 0.220]$
contains $0.178$. The defensible claim is that the ten-dimensional
parameterization \emph{reaches the top cluster}, a $+75\%$ improvement over its
own untuned starting point, not that it is uniquely best.

\begin{figure}[t]
  \centering
  \includegraphics[width=\linewidth]{fig5_memory_comparison.png}
  \caption{Stage~2: twelve memory systems on MultiHopKeyDoor (task performance
  with $95\%$ bootstrap CIs, retrieval precision, token efficiency). \Vfour{}
  sits in the top cluster, statistically tied with the strongest baselines.}
  \label{fig:stage2bench}
\end{figure}

\paragraph{Ablation.} Removing each $\thetavec$ component one at a time
(Figure~\ref{fig:ablation}) shows that $\theta_{\mathrm{novel}}$ is the
dominant, load-bearing dimension: zeroing it collapses reward to $0$ (it gates
the entire storage pipeline). Other components carry less, and a
\wgraph{} sweep (Figure~\ref{fig:wgraph}) shows the graph-traversal term carries
\emph{no} measurable load on this task: zeroing it barely changes reward. The
corpus-tuning signature is therefore a three-shift rather than a four-shift; the
graph term behaves as a typed storage scaffold rather than a relational retrieval
engine.

\begin{figure}[t]
  \centering
  \begin{subfigure}{0.49\linewidth}
    \includegraphics[width=\linewidth]{fig_ablation_ranked.png}
    \caption{Per-component degradation ranking.}
    \label{fig:ablation}
  \end{subfigure}\hfill
  \begin{subfigure}{0.49\linewidth}
    \includegraphics[width=\linewidth]{fig_w_graph_ablation.png}
    \caption{\wgraph{} sweep: flat (no load).}
    \label{fig:wgraph}
  \end{subfigure}
  \caption{Stage~2 ablation. $\theta_{\mathrm{novel}}$ is load-bearing; the graph
  term carries no measurable retrieval load.}
\end{figure}

\paragraph{Transfer and sensitivity.} A $\thetavec$ tuned on one environment
transfers strongly to a structurally similar one (Goal-Room $0.69$, HardKeyDoor
$0.17$) but fails completely on a larger out-of-distribution task
(MegaQuestRoom $0.00$), even though its retrieval there is high-precision ---
the failure is task scale, not memory quality (Figure~\ref{fig:transfer}). The
2-D sensitivity landscape (Figure~\ref{fig:sensitivity}) shows a \emph{sharp
peak}: reward collapses once $\theta_{\mathrm{novel}}$ drops below $\sim$0.3, so
the optimizer must find a fairly precise value. A neural-controller variant
(a $\sim$2{,}000-parameter MLP tuned by the same CMA-ES budget) matches but does
not clearly exceed the ten-scalar \Vfour{}; details are in
Appendix~\ref{app:tables}. The Pareto view (Figure~\ref{fig:pareto}) confirms
\Vfour{} is on the reward-vs-cost frontier.

\begin{figure}[t]
  \centering
  \begin{subfigure}{0.49\linewidth}
    \includegraphics[width=\linewidth]{fig_transfer_annotated.png}
    \caption{Zero-shot $\thetavec$-transfer.}
    \label{fig:transfer}
  \end{subfigure}\hfill
  \begin{subfigure}{0.49\linewidth}
    \includegraphics[width=\linewidth]{fig_sensitivity_annotated.png}
    \caption{Sensitivity landscape (sharp peak).}
    \label{fig:sensitivity}
  \end{subfigure}
  \caption{Stage~2 transfer and sensitivity.}
\end{figure}

\begin{figure}[t]
  \centering
  \includegraphics[width=0.8\linewidth]{fig11_pareto.png}
  \caption{Stage~2: reward versus token cost. \Vfour{} sits on the efficient
  frontier (high reward, low retrieval cost).}
  \label{fig:pareto}
\end{figure}

\subsection{What the three Stage-2 findings jointly imply about the
parameterization}
\label{sec:stage2structure}
The benchmark, ablation, transfer, and sensitivity results in
Section~\ref{sec:stage2} are usually read one at a time, but taken together they
triangulate a single structural fact about how the ten-dimensional $\thetavec$
is organized. Three of those results point the same way.

\paragraph{Within-family-but-not-across-family transfer localizes the useful
variance.} The transfer result (Figure~\ref{fig:transfer}) is not a simple
``transfer is imperfect'' story. A $\thetavec$ tuned on one environment transfers
\emph{strongly} to a structurally similar one (Goal-Room $0.69$, HardKeyDoor
$0.17$) yet collapses to $0.00$ on the larger out-of-distribution task --- and
crucially it collapses \emph{while retaining high retrieval precision there}. If
the learned vector had memorized surface idiosyncrasies of its training
environment, transfer would degrade smoothly with distance; instead it holds
within a task family and then fails discontinuously at a scale boundary. That
dissociation --- good retrieval, zero reward --- says the failure is not in the
memory policy $\thetavec$ encodes but in the task's answering demands outstripping
what any fixed-$k$ retrieval can supply. The operative variance that $\thetavec$
captures is therefore \emph{family-level structure} (how to store and rank), not
instance-level detail, which is exactly the property a parameterization needs if
it is to be predicted rather than re-tuned per task.

\paragraph{The flat \wgraph{} ablation identifies an inert coordinate.} The
\wgraph{} sweep (Figure~\ref{fig:wgraph}) is flat: zeroing the graph-traversal
term barely changes reward, in sharp contrast to $\theta_{\mathrm{novel}}$, whose
removal collapses reward to $0$ (Figure~\ref{fig:ablation}). Read alongside the
ablation ranking, this means the ten nominal dimensions are not ten
\emph{effective} dimensions: at least one axis (\wgraph{}) is a flat direction of
the objective on this task, a typed storage scaffold rather than a load-bearing
retrieval knob. That the term nonetheless appears to shift under corpus tuning
(the $0\rightarrow1.6$ move discussed in Section~\ref{sec:stage2}) without
changing accuracy is the same observation from the optimizer's side: it drifts
along a direction the loss does not penalize. An inert coordinate is not a defect
to be tuned away --- it is the reason the effective search space is smaller than
its nominal dimension, which is what makes the ten-scalar controller reach the
top cluster on a modest CMA-ES budget and what makes the neural variant unable to
clearly beat it: there is little residual signal left in the extra capacity.

\paragraph{The anisotropic objective ties the two together.} The sensitivity
landscape (Figure~\ref{fig:sensitivity}) shows a sharp peak: reward collapses
once $\theta_{\mathrm{novel}}$ drops below $\sim$$0.3$. Set against the flat
\wgraph{} direction, the objective is strongly \emph{anisotropic} --- knife-edged
along $\theta_{\mathrm{novel}}$, nearly flat along \wgraph{} --- i.e.\ a
ridge-like surface rather than an isotropic bowl. This single geometric picture
explains all three earlier observations at once. The narrow, load-bearing axes
(dominated by $\theta_{\mathrm{novel}}$) are what must be gotten right and what
transfers within a family, because task families share those axes; the flat axes
(\wgraph{}) are what neither transfer nor ablation can detect a signal on,
because moving along them costs nothing. It also reconciles the top-cluster
\emph{tie} in Figure~\ref{fig:stage2bench} with the necessity of tuning: only a
few directions carry the reward, so a well-placed ten-scalar vector reaches the
frontier, but those few directions are sharp enough (the $\sim$$0.3$ cliff) that
the fixed default falls off the ridge --- which is precisely why the untuned
starting point sits $+75\%$ below the tuned optimum. The parameterization, in
short, is low-effective-dimension and anisotropic: a small number of sharp,
family-shared axes doing the work, wrapped in several inert ones. This is the
Stage-2 fact that Section~\ref{sec:thetapredict} later exploits, and it is why a
naive parameter-distance check there is uninformative --- distance along the flat
directions is retrieval-irrelevant.

\section{Stage 3 --- six real long-context benchmarks}
\label{sec:stage3}
Each benchmark is ingested document-by-document into the \Vt{} memory; questions
are answered \emph{online} and \emph{batch} and scored by the cross-vendor
Claude judge (Section~\ref{sec:eval}). $\thetavec$ is re-tuned per benchmark on
the corpus-cumulative \recallk{} objective.

\paragraph{Headline: end-of-corpus accuracy lift.}
Table~\ref{tab:stage3} reports the batch (end-of-corpus) judge means under the
canonical (grid-tuned) $\thetavec$ and the corpus-tuned $\thetavec$, with the
held-out significance verdict.

Of the three \emph{domain-coherent} benchmarks (where the hypothesis predicts a
lift), two --- FinanceBench and CUAD --- show a held-out lift that survives Holm
correction; QASPER does not. HotpotQA and LongMemEval are pre-registered
\emph{domain-incoherent controls} (no lift predicted), and NarrativeQA's tuning
objective is undefined by construction. We therefore lead with the count that the
design supports: \textbf{two of three coherent benchmarks show a significant
end-of-corpus lift}.

\begin{table}[t]
\centering
\caption{Stage~3 end-of-corpus (batch) judge means with bootstrap 95\% CIs,
grouped by benchmark type. Lift = corpus-tuned $-$ canonical $\thetavec$.
Canonical and corpus-tuned CIs are non-overlapping for FB, CUAD, QASPER and HQA;
they touch for LME. ``Survives'' = held-out lift significant after Holm.}
\label{tab:stage3}
\small
\begin{tabular}{l r c c c l}
\toprule
Benchmark & $n$ & canonical [95\% CI] & corpus-tuned [95\% CI] & lift & held-out \\
\midrule
\multicolumn{6}{l}{\textit{Confirmatory (domain-coherent)}}\\
\FB{}  & 150 & 0.243 [0.18,\,0.31] & 0.645 [0.57,\,0.72] & $+0.402$ & \textbf{survives}, $p_{\mathrm{holm}}<10^{-4}$ \\
CUAD   & 644 & 0.028 [0.02,\,0.04] & 0.172 [0.15,\,0.20] & $+0.144$ & \textbf{survives} ($+0.135$) \\
QASPER & 94  & 0.250 [0.18,\,0.32] & 0.415 [0.34,\,0.50] & $+0.165$ & \emph{n.s.} after Holm \\
\midrule
\multicolumn{6}{l}{\textit{Controls (domain-incoherent, pre-registered)}}\\
\HQA{} & 100 & 0.215 [0.14,\,0.29] & 0.755 [0.67,\,0.83] & $+0.540$ & (control) \\
\LME{} & 100 & 0.165 [0.10,\,0.25] & 0.330 [0.24,\,0.42] & $+0.165$ & (control) \\
\midrule
\multicolumn{6}{l}{\textit{Undefined objective}}\\
\NQA{} & 10  & 0.400 [0.15,\,0.68] & 0.400 [0.15,\,0.68] & $0.000$ & undefined \\
\bottomrule
\end{tabular}
\end{table}

The mechanism is consistent across benchmarks: corpus tuning drives \wrec{}
toward $0$ and \wembed{} up, producing a memory that does \emph{not} depend on
the queried document being recent --- exactly what end-of-corpus QA needs. In
\emph{online} mode (fresh document) the recency-heavy canonical $\thetavec$ is
competitive or better, so the lift is genuinely a batch-mode effect
(Figure~\ref{fig:stage3lift}).

\begin{figure}[t]
  \centering
  \includegraphics[width=0.92\linewidth]{fig_stage3_lift.png}
  \caption{Stage~3 end-of-corpus (batch) judge means, canonical vs.\
  corpus-tuned $\thetavec$, across the six benchmarks. Of the three
  domain-coherent benchmarks, the held-out lift survives Holm correction on two
  (FinanceBench, CUAD); QASPER is not significant. HotpotQA and LongMemEval are
  pre-registered domain-incoherent controls and NarrativeQA is
  undefined-by-construction.}
  \label{fig:stage3lift}
\end{figure}

\paragraph{Per-benchmark reading.} The benchmarks do not all
tell the same story, and we report them as tiered evidence rather than a uniform
win. \textbf{FinanceBench} is the cleanest result: a held-out batch lift of
$+0.335$ with $p_{\mathrm{holm}}<10^{-4}$, on questions the CMA-ES tuner never
saw. \textbf{CUAD} survives held-out ($+0.135$) and the lift holds at the
50-contract scale (below). \textbf{HotpotQA} and \textbf{LongMemEval} were
pre-registered as domain-incoherent \emph{controls}; their batch lifts
($+0.540$, $+0.165$ at $n{=}100$) are reported as what the controls show, not as
clean wins. \textbf{QASPER}'s pooled lift ($+0.165$) does \emph{not} survive Holm
correction and is reported as non-significant. \textbf{NarrativeQA} has no
paragraph-level gold, so the tuning objective is undefined and both
configurations score identically by construction. In summary,
\emph{two of three domain-coherent benchmarks (FinanceBench, CUAD) show a
significant held-out lift}; the two domain-incoherent controls also lift (a point
we return to in the Discussion), and NarrativeQA is null by construction --- one
consistent mechanism throughout.

\paragraph{The corpus-tuning signature.} The lift has a clear mechanistic
explanation in how $\thetavec$ shifts under corpus-cumulative tuning. On
FinanceBench the recency weight collapses (\wrec{}: $3.78 \rightarrow 0.003$),
the embedding weight rises (\wembed{}: $1.08 \rightarrow 2.63$), and the storage
threshold tightens ($\theta_{\mathrm{store}}$: $0.29 \rightarrow 0.01$): the
tuned memory stores selectively and retrieves by \emph{semantic match}, not
recency. This is precisely what end-of-corpus QA requires --- the answer was
seen long ago, so a recency-driven retriever fetches the wrong (recent) document.
The same direction of shift recurs across the benchmarks that improve, which is
why we describe it as a \emph{signature} rather than a per-benchmark accident. A
fourth apparent shift --- the graph weight rising
(\wgraph{}: $0 \rightarrow 1.6$) --- is \emph{not} load-bearing: the \wgraph{}
ablation (Section~\ref{sec:stage2}) shows zeroing it barely changes accuracy, so
the operative signature is a three-component shift (down-weight recency,
up-weight embedding, tighten storage), with the graph term inert.

\paragraph{Construct validity: does the tuning objective predict answer quality?}
CMA-ES optimizes recall@8 of the gold evidence, but the headline metric is the
LLM judge score; these need not coincide. We test the link directly by joining
per-question recall@8 with the judge score across all corpus-mode batch cells
($n=2{,}170$ questions, five benchmarks with gold evidence). Retrieval recall is
a strong predictor of answer quality: pooled Spearman $\rho = 0.69$ (per
benchmark $0.39$--$0.84$), and questions for which the gold evidence \emph{is}
retrieved are judged correct far more often than those where it is missed (mean
judge $0.62$ vs.\ $0.10$, Figure~\ref{fig:recalljudge}). The causal chain behind
the headline is therefore intact: corpus tuning raises recall \emph{and} judge
score together on every benchmark (e.g.\ FinanceBench recall $0.22\rightarrow0.97$, judge
$0.243\rightarrow0.645$; HotpotQA recall $0.11\rightarrow1.00$, judge $0.215\rightarrow0.755$). QASPER is
the weakest link ($\rho=0.39$), consistent with its being the only coherent
benchmark whose lift does not survive Holm correction.

\begin{figure}[t]
  \centering
  \includegraphics[width=0.7\linewidth]{fig_recall_vs_judge.png}
  \caption{Construct validity: mean LLM judge score as a function of retrieval
  recall@8 (the CMA-ES tuning objective), pooled over $2{,}170$ corpus-mode
  questions. When the gold evidence is retrieved (recall $=1$) answers score
  $0.62$ on average; when it is missed (recall $=0$), $0.10$. Pooled Spearman
  $\rho=0.69$.}
  \label{fig:recalljudge}
\end{figure}

\paragraph{CUAD scaled from 10 to 50 contracts.} To check that the CUAD lift is
not an artifact of a tiny corpus, the headline pair was re-run and re-judged
one-by-one at 50 contracts ($n{=}644$ questions per cell, up from $132$). The
batch lift \emph{holds}: canonical $0.028 \rightarrow$ corpus-tuned $0.172$
($+0.144$, versus $+0.161$ at 10 contracts). Canonical batch stays near-zero
because, with 50 contracts in memory, its recency-heavy retrieval pulls
clauses from the wrong contract; the corpus-tuned $\thetavec$ (with \wrec{}
$\approx 0$) recovers by matching on semantic similarity. The online-vs-batch
contrast is the sharpest of any benchmark
(Table~\ref{tab:cuadonline}): every configuration is far better online (the
right contract is fresh) than batch, and corpus tuning closes most of that gap.

\begin{table}[t]
\centering
\caption{CUAD online-vs-batch judge means at 50 contracts ($n{=}644$). The
online--batch gap is the largest of any benchmark; corpus tuning narrows it.}
\label{tab:cuadonline}
\begin{tabular}{lccc}
\toprule
Configuration & Online & Batch & Gap \\
\midrule
\Vt{} canonical     & 0.295 & 0.028 & $+0.267$ \\
\Vt{} corpus-tuned  & 0.390 & 0.172 & $+0.218$ \\
\bottomrule
\end{tabular}
\end{table}

\begin{figure}[t]
  \centering
  \includegraphics[width=0.7\linewidth]{fig_stage3_online_batch.png}
  \caption{CUAD online vs.\ batch judge means at 50 contracts. Both
  configurations are far stronger online (the source contract is recent); the
  recency-heavy canonical $\thetavec$ collapses at end-of-corpus, while the
  corpus-tuned $\thetavec$ recovers most of the gap.}
  \label{fig:onlinebatch}
\end{figure}

\paragraph{Full-corpus scaling: the lift holds at 510 contracts and 281 papers.}
The headline numbers above were established at reduced corpus scopes --- CUAD at
$50$ contracts, the fair baselines at a $94$-question QASPER subset --- to keep
the one-by-one judging tractable. Because the corpus-cumulative harness holds the
\emph{entire} corpus in memory for the end-of-corpus batch regardless of scope, we
can read the already-judged batch cells at the \emph{complete} corpus --- all
$510$ CUAD contracts ($n{=}6702$ questions) and all $281$ QASPER papers
($n{=}1005$) --- and ask whether the conclusions survive
(Table~\ref{tab:fullcorpus}). They do, and they sharpen. On CUAD the
corpus-tuning lift is significant at every scale we measured and \emph{monotone}
in the expected direction: $+0.161$ at $10$ contracts, $+0.144$ at $50$, and
$+0.107$ at the full $510$ ($0.043 \rightarrow 0.150$, paired Wilcoxon
$p<10^{-12}$, $n{=}6702$). The lift attenuates as the corpus grows --- more
distractor contracts dilute retrieval --- but does not vanish, and the tuned
$\thetavec$ stays clear of both the canonical start and the grid-tuned variant;
the cross-\emph{method} comparison on this surface-lexical corpus (where a tuned
sparse retriever still leads) is the separate question of
Tables~\ref{tab:fairbaseline} and~\ref{tab:head2head}. On QASPER at its full
$281$-paper corpus the learned memory is \emph{tied for best} among all seven
configurations tested: corpus-tuned $\Vt{}$ ($0.351$) beats its own canonical
start by $+0.218$, the full-context dump-all baseline by $+0.240$, stock Okapi
BM25 by $+0.058$ ($p<10^{-4}$), and the tuned attention retriever by $+0.023$
($p{=}0.033$), every margin paired-Wilcoxon significant after Holm correction ---
and it \emph{ties} the one remaining comparator, a fairly-tuned Okapi BM25 run at
the same scale ($0.367$ vs.\ $0.351$; $+0.016$ for BM25, $p{=}0.30$, n.s.), which
is exactly the statistical dead-heat the reduced-scope fair baseline reported
($n{=}94$; Table~\ref{tab:fairbaseline}). The full-corpus view therefore
\emph{strengthens} the reduced-scope headlines: the corpus-tuning effect is not an
artifact of a small evaluation corpus, and on this semantically-matched
scientific-paper domain the learned policy is tied for best at the complete
corpus --- matching a purpose-tuned lexical retriever and beating every other
configuration.

\begin{table}[t]
\centering
\caption{Full-corpus scaling. Judged end-of-corpus batch accuracy (bootstrap
$95\%$ CI) at the \emph{complete} corpus, from the already-judged batch cells
(mode batch, all questions asked once with the full corpus in memory). Deltas and
$p$-values are versus \Vt{} corpus-tuned (paired Wilcoxon, Holm-corrected within
benchmark). The CUAD block reports the learned-memory configurations (the
corpus-tuning lift); the cross-method fairness comparison on CUAD is in
Tables~\ref{tab:fairbaseline}--\ref{tab:head2head}.}
\label{tab:fullcorpus}
\footnotesize
\begin{tabular}{lccl}
\toprule
Configuration & Judged accuracy & $\Delta$ vs.\ tuned & $p_{\mathrm{Holm}}$ \\
\midrule
\multicolumn{4}{l}{\emph{CUAD --- full 510 contracts ($n{=}6702$)}} \\
\Vt{} corpus-tuned & $\mathbf{0.150}$ {\scriptsize[0.143, 0.156]} & --- & --- \\
\Vt{} tuned (grid) & $0.129$ {\scriptsize[0.124, 0.135]} & $-0.020$ & $<10^{-6}$ \\
\Vt{} canonical    & $0.043$ {\scriptsize[0.038, 0.047]} & $-0.107$ & $<10^{-6}$ \\
\midrule
\multicolumn{4}{l}{\emph{QASPER --- full 281 papers ($n{=}1005$)}} \\
Okapi BM25 (tuned) & $0.367$ {\scriptsize[0.341, 0.393]} & $+0.016$ & $0.30$ (tie) \\
\Vt{} corpus-tuned & $0.351$ {\scriptsize[0.327, 0.377]} & --- & --- \\
Attention (tuned)  & $0.328$ {\scriptsize[0.302, 0.354]} & $-0.023$ & $0.033$ \\
Okapi BM25 (stock) & $0.294$ {\scriptsize[0.269, 0.318]} & $-0.058$ & $<10^{-4}$ \\
\Vt{} canonical    & $0.133$ {\scriptsize[0.115, 0.152]} & $-0.218$ & $<10^{-6}$ \\
dump-all           & $0.111$ {\scriptsize[0.095, 0.128]} & $-0.240$ & $<10^{-6}$ \\
\Vt{} tuned (grid) & $0.091$ {\scriptsize[0.073, 0.109]} & $-0.260$ & $<10^{-6}$ \\
\bottomrule
\end{tabular}
\end{table}

\paragraph{The dump-all baseline: accuracy parity, not a cliff.} A full-context
``dump-all'' baseline concatenates the entire accumulated corpus into the prompt,
truncated only when it exceeds the context window. On \FB{} it scores $0.689$
online / $0.607$ batch --- \emph{statistically
indistinguishable} from corpus-tuned ($0.678$ / $0.645$; confidence intervals
overlap heavily) --- but at roughly $18\times$ the token cost (\$2.68 vs.\ \$0.15
per 150 questions). The conclusion is that the value of learned selective
memory at this corpus scale is \textbf{efficiency and scalability}, not a raw
accuracy advantage: it matches dump-all's accuracy at a fraction of the cost and,
unlike dump-all, does not break when the corpus exceeds the context window.

\begin{figure}[t]
  \centering
  \includegraphics[width=0.78\linewidth]{fig_stage3_cost.png}
  \caption{FinanceBench accuracy vs.\ cost. Dump-all is statistically tied with
  corpus-tuned selective memory on accuracy but $\approx 18\times$ more
  expensive; corpus-tuned sits in the cheap, accurate corner.}
  \label{fig:cost}
\end{figure}

\paragraph{Scalability: dump-all structurally breaks; selective memory does not.}
The efficiency argument is not merely about cost --- past a corpus size, dump-all
\emph{cannot run at all}, while selective retrieval is unaffected. We measure this
directly on CUAD (Figure~\ref{fig:scalability}). Each contract averages
$\sim$10.9K tokens, so the concatenated dump-all prompt grows $O(N)$ and crosses
\texttt{gpt-4o-mini}'s 128K-token context window at just $N \approx 11$
contracts; at the full $N{=}510$ it would require \textbf{5.55M tokens}
($\approx 43\times$ the limit), forcing truncation that discards most of the
evidence. Selective retrieval ($k{=}8$) holds its prompt at a measured mean of
\textbf{704 tokens} regardless of $N$. This is the structural case for learned
selective memory: beyond a few documents, retrieving the right $k$ items is not an
optimization but a necessity, and the only open question is how good the retrieval
policy is --- which is exactly what $\thetavec$ tunes.

\begin{figure}[t]
  \centering
  \includegraphics[width=0.85\linewidth]{fig_stage3_scalability.png}
  \caption{Answer-prompt size versus corpus size on CUAD (log scale). The
  full-context dump-all prompt grows linearly and exceeds the answerer's
  128K-token context at $N\approx 11$ contracts (reaching $\approx 43\times$ the
  limit at $N{=}510$); selective retrieval stays flat at $\sim$704 tokens. Beyond
  $N\approx 11$, dump-all must truncate and lose evidence; selective memory is
  unaffected.}
  \label{fig:scalability}
\end{figure}

\paragraph{Fair retrieval baselines: competitive, not dominant.} The dump-all
comparison controls for context but not for \emph{retrieval tuning}: a skeptic
could object that \Vt{} only wins because the classical baselines were left at
their defaults. We therefore tuned the strongest lexical baseline (Okapi BM25,
grid-searching $k_1$ and $b$) on the \emph{same} corpus-recall objective and
budget as $\thetavec$, then judged its answers one-by-one at matched document
scope. Table~\ref{tab:fairbaseline} reports the head-to-head, and the honest
verdict is that learned selective memory is \textbf{competitive but
benchmark-dependent}. It beats fairly-tuned BM25 on \FB{}
($+0.132$, $p_{\mathrm{Holm}}{<}10^{-3}$), is statistically tied on QASPER
($-0.011$, n.s.), and is \emph{beaten} on CUAD ($-0.131$,
$p_{\mathrm{Holm}}{<}10^{-13}$, with non-overlapping confidence intervals). The
CUAD reversal is not noise: it is consistent with the $\thetavec$-predictability
analysis below (Section~\ref{sec:thetapredict}), where CUAD is the most
idiosyncratic corpus and recovers the least tuning lift. Pure clause extraction
rewards exact lexical matching --- which tuned BM25 performs directly and the
graph, entity, and temporal machinery of \Vt{} does not. We therefore make no
claim that learned memory dominates strong retrieval baselines in general: its
accuracy advantage is specific to corpora where relevance is semantic or
relational (financial and narrative QA) rather than lexical, and its
\emph{domain-general} value remains the efficiency and scalability argument
established above.

\begin{table}[t]
  \centering
  \caption{Fair-baseline head-to-head: corpus-tuned \Vt{} versus Okapi BM25 tuned
  on the \emph{same} corpus-recall budget, both judged one-by-one by Claude at
  matched document scope (batch mode, seed~42). Cells are mean judge score with a
  bootstrap $95\%$ CI; $\Delta$ is learned minus baseline; the verdict is from a
  paired Wilcoxon signed-rank test, Holm-corrected across the baseline family.
  Learned memory wins on one benchmark, ties on one, and loses on one.}
  \label{tab:fairbaseline}
  \footnotesize
  \begin{tabular}{lccrl}
    \toprule
    Benchmark & \Vt{} corpus-tuned & BM25 (tuned) & $\Delta$ & Verdict \\
    \midrule
    \FB{} ($n{=}150$) & $\mathbf{0.645}$ {\scriptsize[0.57, 0.72]} & $0.513$ {\scriptsize[0.45, 0.59]} & $+0.132$ & learned wins \\
    QASPER ($n{=}94$) & $0.415$ {\scriptsize[0.33, 0.49]} & $0.426$ {\scriptsize[0.35, 0.50]} & $-0.011$ & tie (n.s.) \\
    CUAD ($n{=}644$) & $0.172$ {\scriptsize[0.15, 0.20]} & $\mathbf{0.303}$ {\scriptsize[0.27, 0.33]} & $-0.131$ & baseline wins \\
    \bottomrule
  \end{tabular}
\end{table}

\paragraph{Head-to-head against HippoRAG and MemGPT/Letta.} The fair-baseline
comparison pits learned memory against tuned classical retrievers; a stronger test
is against two prominent \emph{memory systems}. We compared \Vt{} corpus-tuned
head-to-head with faithful reimplementations of HippoRAG's graph-topology
retrieval \citep{gutierrez2024hipporag} and MemGPT/Letta's tiered archival
retrieval \citep{packer2023memgpt}, both run under the identical corpus-cumulative
harness, the same fixed \texttt{gpt-4o-mini} answerer, the same MiniLM encoder,
and the same Claude one-by-one judge, so that only the \emph{retrieval} component
varies (Table~\ref{tab:head2head}).

The result mirrors --- and sharpens --- the benchmark-dependence of the fair
baselines. On FinanceBench the learned memory ($0.645$) is statistically
\emph{tied} with both systems: numerically ahead of HippoRAG ($0.563$; $+0.082$,
$p_{\mathrm{Holm}}{=}0.06$) and even with Letta ($0.607$; $+0.038$, n.s.), all
intervals overlapping. On QASPER --- also semantically matched, but
information-seeking over scientific papers rather than financial filings --- the
learned memory instead pulls \emph{clearly ahead} of both systems: \Vt{}
($0.351$) beats HippoRAG ($0.286$; $+0.065$, $p_{\mathrm{Holm}}{<}10^{-6}$) and
Letta ($0.319$; $+0.032$, $p_{\mathrm{Holm}}{=}0.001$), both paired-Wilcoxon
significant, and edges the matched-scope tuned-attention ($0.328$) and stock-BM25
($0.294$) references. The judged rationales make the mechanism concrete:
HippoRAG's entity-seeded graph, running on the rule-based extractor rather than an
LLM OpenIE module, repeatedly surfaces the wrong passage and abstains (``the
method is not specified''), and every such refusal scores $0$, whereas the tuned
semantic memory retrieves the short-span evidence QASPER rewards. (The
\emph{fairly-tuned} BM25 retriever, run at this same $n{=}1005$ full-corpus scope,
is a statistical \emph{tie} with \Vt{} --- $0.367$ vs.\ $0.351$, $p{=}0.30$, so
BM25 is nominally ahead but indistinguishable --- the same dead-heat the $n{=}94$
fair baseline reported (Table~\ref{tab:fairbaseline}); the learned memory beats
the two published memory systems here but does not out-retrieve a purpose-tuned
lexical index on this corpus.) On CUAD the ranking \emph{inverts} and the gap turns
against the learned memory: HippoRAG ($0.222$) \textbf{beats} \Vt{} ($0.172$) by a
significant margin ($p_{\mathrm{Holm}}{<}0.01$), \Vt{} ties Letta ($0.146$;
$+0.026$, n.s.), and both trail the tuned lexical retrievers ($0.303$ BM25,
$0.320$ attention). The pattern is coherent with the task-structure argument of
Section~\ref{sec:whybm25}: legal clause extraction rewards \emph{lexical and
entity} matching, so HippoRAG's entity-seeded graph retrieval and BM25's term
overlap both outrank a semantic memory, whereas Letta's embedding-plus-recency
retrieval --- like \Vt{}'s semantics --- has the least to offer, and its recency
core actively surfaces wrong-contract passages. No single retrieval mechanism
dominates across the three corpora, and the split is clean along the
semantic/lexical axis: the learned memory is \emph{best} on QASPER, \emph{tied}
on FinanceBench, and \emph{loses} on CUAD, while HippoRAG is exactly the mirror
image --- weakest on both semantic corpora and strongest on the lexical one. The
learned policy wins or ties where relevance is semantic and forfeits its edge
where relevance is surface-lexical, exactly the regime in which a purpose-built
lexical or entity retriever should win --- which is the empirical premise of a
per-task optimizer: if one fixed mechanism sufficed everywhere there would be
nothing to re-point.

\begin{table}[t]
  \centering
  \caption{External-system head-to-head: mean judged accuracy (bootstrap $95\%$
  CI) under the identical answerer, encoder, and one-by-one Claude judge; only
  retrieval varies. HippoRAG and MemGPT/Letta are faithful reimplementations of
  the published \emph{retrieval mechanisms} (disclosed divergences below); the
  tuned lexical retrievers are shown for reference. No mechanism dominates all
  three corpora: the learned memory wins or ties on the two semantic corpora and
  loses on the lexical one. All QASPER cells are at the full $n{=}1005$ scope; the
  learned memory ties the fairly-tuned BM25 reference there ($0.351$ vs.\ $0.367$,
  $p{=}0.30$), mirroring the $n{=}94$ fair baseline (Table~\ref{tab:fairbaseline}).}
  \label{tab:head2head}
  \footnotesize
  \begin{tabular}{lccc}
    \toprule
    Retrieval mechanism & FinanceBench ($n{=}150$) & QASPER ($n{=}1005$) & CUAD ($n{=}644$) \\
    \midrule
    \Vt{} corpus-tuned (learned)   & $0.645$ {\scriptsize[0.57, 0.72]} & $\mathbf{0.351}$ {\scriptsize[0.33, 0.38]} & $0.172$ {\scriptsize[0.15, 0.20]} \\
    HippoRAG (reimpl.)             & $0.563$ {\scriptsize[0.49, 0.64]} & $0.286$ {\scriptsize[0.26, 0.31]} & $\mathbf{0.222}$ {\scriptsize[0.19, 0.25]} \\
    MemGPT/Letta (reimpl.)         & $0.607$ {\scriptsize[0.53, 0.68]} & $0.319$ {\scriptsize[0.29, 0.34]} & $0.146$ {\scriptsize[0.12, 0.17]} \\
    \midrule
    Okapi BM25 (tuned, ref.)       & $0.513$ {\scriptsize[0.45, 0.59]} & $\mathbf{0.367}$ {\scriptsize[0.34, 0.39]} & $0.303$ {\scriptsize[0.27, 0.33]} \\
    Attention (tuned, ref.)        & $\mathbf{0.657}$ {\scriptsize[0.59, 0.73]} & $0.328$ {\scriptsize[0.30, 0.35]} & $\mathbf{0.320}$ {\scriptsize[0.25, 0.39]} \\
    \bottomrule
  \end{tabular}
\end{table}

Both comparison systems are faithful in-repo reimplementations of the published
\emph{retrieval mechanisms}, not wrappers around the authors' code, and each
carries a disclosed divergence. The HippoRAG reimplementation preserves the
system's defining mechanics --- a schemaless phrase/passage knowledge graph,
embedding-synonymy edges above cosine $\tau{=}0.8$, node specificity
$s_i{=}1/|P_i|$ as an IDF surrogate, Personalized PageRank (damping $0.5$) seeded
on specificity-weighted query entities, and passage ranking by summed PPR mass ---
but replaces the paper's LLM-based OpenIE triple extraction with the thesis's
rule-based entity extractor and intra-passage co-occurrence edges, making zero LLM
calls at ingest or query (which lowers its recall ceiling relative to the
published system, whose own error analysis attributes most error to NER and OpenIE
quality). The Letta reimplementation reproduces its tiered \emph{read} surface ---
an always-on recency core unioned with an embedding-similarity archival search ---
but not MemGPT's defining feature, the LLM agent that self-manages its memory; it
therefore measures Letta's archival retrieval rather than its agentic memory
management. Any measured difference thus reflects retrieval topology and ranking,
not the generator, the encoder, or the LLM-driven stages the originals additionally
employ.

\section{Is \texorpdfstring{$\thetavec$}{theta} predictable from corpus statistics?}
\label{sec:thetapredict}
Section~\ref{ch:results} establishes that the optimal $\thetavec$ is
task-dependent; a natural follow-up is whether it is nonetheless
\emph{predictable} from cheap, LLM-free corpus statistics --- which would convert
``re-tune per task'' into a one-shot prediction. We test this directly. We tune
$\thetavec$ (CMA-ES on \recallk{}, no LLM) on $40$ random $30$-document slices,
eight per benchmark, recording for each slice a nine-feature descriptor (corpus
size, document and passage length, lexical diversity, question length, gold
concentration, gold relative position, entity density) alongside its tuned
$\thetavec$. We then run a \emph{leave-one-benchmark-out} test: a $k$-NN
predictor ($k{=}3$) maps the held-out slice's descriptor to a $\thetavec$ using
only the \emph{other} benchmarks' slices, and we evaluate \recallk{} under that
predicted $\thetavec$ on the held-out slice.

The predicted $\thetavec$ --- never tuned on the target benchmark --- recovers a
mean $\mathbf{0.80}$ of the per-task tuning lift (bootstrap $95\%$ CI
$[0.71, 0.87]$): mean recall climbs canonical $0.092 \rightarrow$ predicted
$0.643 \rightarrow$ tuned $0.704$ (Figure~\ref{fig:thetatransfer}). Recovery is
near-total on \HQA{} ($0.99$), \LME{} ($0.96$) and \FB{} ($0.93$), partial on
QASPER ($0.74$), and lowest on \textbf{CUAD} ($0.37$) --- the most idiosyncratic
corpus, where per-task tuning remains necessary. Two honest qualifications: (i)
this is recovery of the \emph{retrieval objective} (\recallk{}), not of judged
answer quality, though the two are correlated ($\rho{=}0.69$,
Figure~\ref{fig:recalljudge}); and (ii) the high average is inflated by a
pathological canonical default (mean canonical recall $0.092$), so the genuine
per-task residual is the remaining $\sim$$20\%$ overall and $\sim$$63\%$ on CUAD.
Notably, a naive pairwise descriptor-distance vs.\ $\thetavec$-distance check is
uninformative ($\rho{=}-0.38$): two $\thetavec$ can be far apart in parameter
space yet retrieval-equivalent (the objective is flat along inert directions such
as \wgraph{}), so recall recovery --- not parameter proximity --- is the valid
test.

\begin{figure}[t]
  \centering
  \includegraphics[width=0.78\linewidth]{fig_theta_transfer.png}
  \caption{Leave-one-benchmark-out: fraction of the per-task tuning lift in
  \recallk{} recovered by a $\thetavec$ predicted (k-NN over corpus descriptors)
  from the \emph{other} benchmarks. Mean $0.80$ (dashed; $95\%$ CI
  $[0.71,0.87]$). Transfer is near-total except on CUAD ($0.37$), the hardest,
  most idiosyncratic corpus.}
  \label{fig:thetatransfer}
\end{figure}
